\documentclass[a4paper,10pt]{article}

\usepackage{latexsym}
\usepackage[empty]{fullpage}
\usepackage{titlesec}
\usepackage{marvosym}
\usepackage[usenames,dvipsnames]{color}
\usepackage{verbatim}
\usepackage{enumitem}
\usepackage[pdftex]{hyperref}
\usepackage{fancyhdr}
\usepackage{mathptmx}
\usepackage{csquotes}

\pagestyle{fancy}
\fancyhf{}
\renewcommand{\headrulewidth}{0pt}
\renewcommand{\footrulewidth}{0pt}

% Adjust margins
\addtolength{\oddsidemargin}{-0.530in}
\addtolength{\evensidemargin}{-0.375in}
\addtolength{\textwidth}{1in}
\addtolength{\topmargin}{-0.5in}
\addtolength{\textheight}{1.1in}

\urlstyle{rm}

\raggedbottom
\raggedright
\setlength{\tabcolsep}{0in}

% Sections formatting
\titleformat{\section}{
  \vspace{-10pt}\scshape\raggedright\large
}{}{0em}{}[\color{black}\titlerule \vspace{-6pt}]

% Custom commands
\newcommand{\resumeItem}[2]{
  \item{
    \textbf{#1}{: #2 \vspace{-2pt}}
  }
}

\newcommand{\resumeSubheading}[4]{
  \vspace{-1pt}\item
    \begin{tabular*}{0.97\textwidth}{l@{\extracolsep{\fill}}r}
      \textbf{#1} & #2 \\
      \textit{#3} & \textit{#4} \\
    \end{tabular*}\vspace{-5pt}
}

\newcommand{\resumeSubItem}[2]{\resumeItem{#1}{#2}\vspace{-3pt}}

\renewcommand{\labelitemii}{$\circ$}

\newcommand{\resumeSubHeadingListStart}{\begin{itemize}[leftmargin=*]}
\newcommand{\resumeSubHeadingListEnd}{\end{itemize}\vspace{-2pt}}
\newcommand{\resumeItemListStart}{\begin{itemize}}
\newcommand{\resumeItemListEnd}{\end{itemize}\vspace{-6pt}}

\begin{document}

% Heading
\begin{tabular*}{\textwidth}{l@{\extracolsep{\fill}}r}
  \textbf{\Large Junyi (Conny) Zhou} & Tel: (404) 663-5601 \\
  Boston, MA 02215 & Email: \href{mailto:junyizhou@hsph.harvard.edu}{junyizhou@hsph.harvard.edu} \\
  \href{https://github.com/JunyiZhou-Conny}{github.com/JunyiZhou-Conny} & \href{https://www.linkedin.com/in/junyi-zhou-270208247}{linkedin.com/in/junyi-zhou-270208247} \\
\end{tabular*}

\section{Education}
\resumeSubHeadingListStart
  \resumeSubheading
  {Harvard University}{Boston, MA}
  {M.S. in Health Data Science}{August 2025 -- March 2027}
  \resumeItemListStart
    \item Program requirements complete December 2026 (commencement March 2027); available for full-time start January 2027
    \item Coursework: Deep Learning, AI in Medicine, Healthcare Machine Learning, Data Science
  \resumeItemListEnd
  \resumeSubheading
  {Emory University, College of Arts and Sciences}{Atlanta, GA}
  {B.S. in Applied Mathematics and Statistics; Minor in Computer Informatics; GPA: 3.925}{August 2021 -- May 2025}
  \resumeItemListStart
    \item Coursework: Probability and Statistics, Machine Learning, NLP, Data Structures, Database Systems
  \resumeItemListEnd
\resumeSubHeadingListEnd

% Technical Skills
\section{Technical Skills}
\resumeSubHeadingListStart
  \resumeItem{Languages}{Python, C++, Java, SQL, R, JavaScript}
  \resumeItem{ML \& Deep Learning}{PyTorch, TensorFlow, Keras, scikit-learn, XGBoost, LightGBM, Hugging Face Transformers}
  \resumeItem{LLM \& Applied AI}{RAG, agentic pipelines, prompt engineering, embeddings \& vector search, OpenAI API, MongoDB Atlas Vector Search}
  \resumeItem{Data \& Compute}{PostgreSQL, MySQL, MongoDB, Spark, Hadoop, SLURM/HPC, multi-GPU job scheduling}
  \resumeItem{Cloud}{AWS (SageMaker, Lambda, EC2, S3, Redshift, Elastic Beanstalk, CloudFront), Snowflake, BigQuery}
  \resumeItem{Backend \& DevOps}{Flask, REST APIs, React, Docker, Auth0, Git}
  \resumeItem{Certifications}{AWS Certified Cloud Practitioner, AWS Certified AI Practitioner, AWS Machine Learning Associate}
\resumeSubHeadingListEnd

% Clinical AI Experience
\section{Clinical AI Experience}
\resumeSubHeadingListStart

  \resumeSubheading
    {Airway Management Simulation Chatbot \href{https://d3g1qw60hz7yw7.cloudfront.net}{\normalfont\small [Live Demo]}}{Atlanta, GA}
    {Full-Stack Engineer, Scrum Master | Emory Pediatric Hospital}{January 2024 -- August 2025}
    \resumeItemListStart
      \item Architected a clinical \textbf{agentic LLM pipeline} pairing \textbf{RAG} retrieval over \textbf{MongoDB Atlas Vector Search} with \textbf{prompt engineering} and context-window management to generate multi-turn training scenarios from \textbf{OpenAI} models.
      \item Built a \textbf{HIPAA-conscious} full-stack platform (\textbf{Python/Flask} REST API, React, MongoDB, \textbf{Auth0}) deployed on \textbf{AWS} Elastic Beanstalk and CloudFront, tracking resident performance metrics.
      \item Led an Agile team of 6, running code reviews and folding clinician feedback into successive releases.
    \resumeItemListEnd

  \resumeSubheading
    {Cross-Species Drug Translation using Optimal Transport \& Generative AI}{Boston, MA}
    {Graduate Researcher | Mooney Lab \& Alvarez-Melis Lab, Wyss Institute}{February 2026 -- Present}
    \resumeItemListStart
      \item Built \textbf{Optimal Transport} models (\textbf{VAEs}, \textbf{unbalanced OT}) predicting animal-to-human trial translation from \textbf{scRNA-seq} atlases, with \textbf{scanpy} pipelines automating cell-type labeling and dataset integration.
      \item Ported a 2018 CPU-only \textbf{TensorFlow} reference implementation to \textbf{PyTorch}, reimplementing the \textbf{Input Convex Neural Networks} behind the optimal-transport dual and enabling GPU training on Harvard's \textbf{FASRC Cannon} cluster.
      \item Benchmarked \textbf{CellOT} against \textbf{scGen} on cross-species perturbation prediction, measuring 0.85--0.90 $R^2$ with the source species in training versus 0.65--0.67 on the held-out-species split.
      \item Established through \textbf{paired single-factor ablations} that gene-space choice reverses with split --- 6,619 orthologs wins in-distribution while 1,000 HVGs wins on every out-of-distribution metric --- ruling out a split-independent default.
    \resumeItemListEnd

\resumeSubHeadingListEnd

% Machine Learning Projects
\section{Machine Learning Projects}
\resumeSubHeadingListStart

  \resumeSubheading
    {Pneumonia Detection from Chest Radiographs \href{https://github.com/JunyiZhou-Conny/BST-261-2026-Pneumonia-Chest-X-Ray-Classification}{\normalfont\small [GitHub]}}{Boston, MA}
    {Harvard BST-261 (Kaggle competition)}{Spring 2026}
    \resumeItemListStart
      \item Diagnosed a 14-point generalization gap on a \textbf{Kaggle} chest-radiograph benchmark --- 99\% validation against 85\% leaderboard --- and traced it to memorization of a 4,700-image training set.
      \item Closed the gap to 96\% with domain-aware augmentation (\textbf{Mixup} $\alpha = 0.2$, RandomResizedCrop, RandomErasing, affine jitter), excluding vertical flips that would render chest radiographs anatomically implausible.
      \item Found that freezing \textbf{DenseNet-121}'s early layers outweighed architecture choice, beating ResNet-50, ConvNeXt-Tiny, EfficientNet-B0, and a three-model ensemble to reach a 0.959 public score.
    \resumeItemListEnd

  \resumeSubheading
    {AlphaFold Protein--Nucleic Acid Interaction Prediction \href{https://github.com/JunyiZhou-Conny/AlphaFold-Bou-Nader-Lab}{\normalfont\small [GitHub]}}{Atlanta, GA}
    {Data Analyst | Bou-Nader Lab, Emory Dept. of Biochemistry}{February 2025 -- August 2025}
    \resumeItemListStart
      \item Engineered a pipeline integrating \textbf{AlphaFold-Multimer} and \textbf{AlphaFold3} to predict protein--nucleic acid interactions from \textbf{MSA} data.
      \item Automated high-throughput inference across GPU nodes with \textbf{SLURM} parallel scheduling, tuning batch size to raise GPU utilization on shared \textbf{HPC} clusters.
      \item Built \textbf{Python} tooling (PyMOL, ChimeraX) ranking structures by \textbf{pLDDT} and \textbf{PAE} score, shipped as a documented CLI so non-technical researchers could run inference and triage results independently.
    \resumeItemListEnd

\resumeSubHeadingListEnd

\end{document}
